\begin{frame}{Research Gaps}
\begin{block}{Summary of Literature Review Observations}
While several assistive technologies exist, existing solutions are constrained: either requiring expensive specialized hardware, lacking deep scene semantics, or failing to address the complexities of unstructured Indian road environments.
\end{block}   
\begin{columns}[T]
  \begin{column}{0.48\textwidth}
    \begin{alertblock}{Gap 1: Expensive Dedicated Hardware}
      \begin{itemize}
        \item Systems like SmartCane, Echo-Vision, and ALVU require custom sensors (LiDAR, ultrasonic modules, smart glasses).
        \item Substantially raises costs and creates adoption barriers.
      \end{itemize}
    \end{alertblock}
    \begin{alertblock}{Gap 2: Limited Scene Semantics}
      \begin{itemize}
        \item Existing aids only output obstacle distance or bounding boxes.
        \item Lack Vision-Language Models (VLM) for rich environmental understanding.
      \end{itemize}
    \end{alertblock}
  \end{column}
  \begin{column}{0.48\textwidth}
    \begin{alertblock}{Gap 3: No End-to-End Smartphone Integration}
      \begin{itemize}
        \item No reviewed system unifies real-time YOLO detection, VLM scene explanation, GPS navigation, OCR, and voice dialogue in one app.
      \end{itemize}
    \end{alertblock}
    \begin{alertblock}{Gap 4: Indian Road Environments}
      \begin{itemize}
        \item Indian roads exhibit mixed traffic, potholes, vendors, and broken pavements that Western benchmarked systems fail to handle.
      \end{itemize}
    \end{alertblock}
  \end{column}
\end{columns}
\end{frame}